\subsection{The cross-modal panel across four open-weight models}
\label{sec:results-crossmodal}

Section~\ref{sec:results-frontier} measured one frontier checkpoint through a relay. One
model, however large, cannot separate the mechanism from a single serving endpoint. We
therefore ran the same six cells on the same three non-CXR modalities for the three local
checkpoints of the main matrix: 1,788 images (600 brain MRI, 600 retinal OCT, 588
dermoscopy), five cells under the v0 contract plus one mid-gray blank under the v1
contract, 32,184 calls. The three cohorts are public, de-identified releases
\citep{cheng2017braintumor,kermany2018identifying,tschandl2018ham10000}; Appendix~\ref{app:crossmodal-data}
gives their source, licence, label mapping, sampling rule and split unit. Images enter at
stored resolution; prompts, label sets, operators
and parser are shared with the CXR arm. The protocol differs from the relay arm in two
ways: the absence of the 1,024-pixel client downscale, and the token budget (1,024
locally; 8,192 on the relay, where hidden reasoning otherwise consumes the whole budget).

146 of the 2,940 dermoscopy cells (5.0\%) exhausted the 1,024-token budget inside
MedGemma-1.5's \texttt{<unused94>thought} channel. A separate repair arm re-ran exactly
those cells at 2,048 tokens (141 recovered) and the remaining five at 4,096 (all five
recovered), for 151 calls and 100\% parse. The repair is reported as a coverage
sensitivity and is not merged into the primary file; Table~\ref{tab:crossmodal4} uses the
frozen 1,024 protocol and reports the parse and completeness rates of every cell. At the
repaired budget, MedGemma-1.5's dermoscopy clean set moves from 95 to 114 distinct vectors
and its label completeness from 0.678 to 0.784, while P1 contrast moves only from 0.549 to
0.563 and P1 noise is unchanged---the primary table is conservative on this axis, not
flattering.

\begin{table}[t]
\centering
\footnotesize
\setlength{\tabcolsep}{3pt}
\caption{Cross-modal panel across four open-weight checkpoints: three run locally, one
served through an OpenAI-compatible relay. ``Distinct'' is the number of different status
vectors on clean images and ``modal'' the share of the most common one. P1$_{\text{c}}$
and P1$_{\text{n}}$ are the share of images whose status vector changes under
contrast-severe and noise-severe relative to clean. $\Delta_{\text{null}}$ is
$D(\text{real},\text{blank})-D(\text{real},\text{real}')$ for the mid-gray v0 blank, with
a 95\% patient-, image- or lesion-clustered bootstrap interval. ``Parse'' and ``Compl.''
are the share of clean cells that parse and the mean share of contract labels named. The
relay rows use 300 images per modality; the local rows use the full manifests. The relay
rows repeat the corresponding rows of Table~\ref{tab:api-null} and are included here for
the four-model comparison. ``Support'' groups Distinct and Modal; ``Sensitivity''
groups P1$_{\text{c}}$ and P1$_{\text{n}}$. For every model the dermoscopy row has the largest Distinct,
P1$_{\text{c}}$ and P1$_{\text{n}}$; $^{\dagger}$ marks a nominal 95\% interval whose cell
does not survive a Bonferroni correction over the 12 $\Delta_{\text{null}}$ estimates.}
\label{tab:crossmodal4}
\begin{tabular}{l l r r r r r l r r}
\toprule
& & & \multicolumn{2}{c}{Support} & \multicolumn{2}{c}{Sensitivity} & & & \\
\cmidrule(lr){4-5}\cmidrule(lr){6-7}
Dataset & Model & $n$ & Distinct & Modal & P1$_{\text{c}}$ & P1$_{\text{n}}$ & $\Delta_{\text{null}}$ [95\% CI] & Parse & Compl. \\
\midrule
Dermoscopy & MedGemma-1.5-4B & 588 & 95 & 0.111 & 0.549 & 0.707 & 0.291 [0.256, 0.327] & 0.889 & 0.678 \\
           & InternVL3.5-8B & 588 & 28 & 0.284 & 0.289 & 0.466 & 0.421 [0.381, 0.461] & 1.000 & 1.000 \\
           & Qwen3-VL-8B & 588 & 12 & 0.488 & 0.184 & 0.355 & 0.504 [0.464, 0.542] & 1.000 & 0.187 \\
           & DeepSeek-V4.1-Flash (relay) & 300 & 94 & 0.117 & 0.417 & 0.540 & 0.155 [0.093, 0.215] & 1.000 & 1.000 \\
\midrule
Brain MRI  & MedGemma-1.5-4B & 600 & 9 & 0.445 & 0.097 & 0.272 & 0.209 [0.136, 0.282] & 1.000 & 1.000 \\
           & InternVL3.5-8B & 600 & 4 & 0.470 & 0.083 & 0.203 & 0.140 [0.068, 0.215] & 1.000 & 1.000 \\
           & Qwen3-VL-8B & 600 & 10 & 0.665 & 0.057 & 0.227 & $-0.059$ [$-0.112$, $-0.003$]$^{\dagger}$ & 1.000 & 1.000 \\
           & DeepSeek-V4.1-Flash (relay) & 300 & 10 & 0.207 & 0.170 & 0.327 & 0.203 [0.122, 0.284] & 1.000 & 1.000 \\
\midrule
Retina OCT & MedGemma-1.5-4B & 600 & \textbf{3} & \textbf{0.928} & 0.087 & 0.087 & $-0.058$ [$-0.080$, $-0.038$] & 1.000 & 0.984 \\
           & InternVL3.5-8B & 600 & 6 & 0.353 & 0.103 & 0.108 & 0.474 [0.428, 0.519] & 1.000 & 1.000 \\
           & Qwen3-VL-8B & 600 & 8 & 0.462 & 0.128 & 0.335 & 0.057 [0.005, 0.105]$^{\dagger}$ & 1.000 & 0.859 \\
           & DeepSeek-V4.1-Flash (relay) & 300 & 16 & 0.307 & 0.187 & 0.207 & 0.390 [0.320, 0.460] & 1.000 & 1.000 \\
\bottomrule
\end{tabular}
\end{table}

\paragraph{The modality ordering is model-invariant.} Dermoscopy is the most responsive of
the three modalities for all four models: P1 contrast is 0.184--0.549 and P1 noise
0.355--0.707, against 0.057--0.170 and 0.087--0.335 on brain MRI and retinal OCT. The
margin over retinal OCT is small for Qwen3-VL (0.184 vs.\ 0.128 and 0.355 vs.\ 0.335),
which is the model whose dermoscopy replies name only 19\% of the contract's labels, but
the ordering itself does not flip in any model--modality pair. This is a within-model
comparison at each model's native geometry. The controlled sweep in
Section~\ref{sec:results-geometry} shows that cross-model P1 rankings can reverse with the
input edge; the P1 columns of Table~\ref{tab:crossmodal4} are native-geometry values, not
geometry-controlled rankings. Distinct status vectors
follow the same pattern: 12--95 on dermoscopy, 4--10 on brain MRI and 3--16 on retinal OCT.
The extreme is MedGemma-1.5 on OCT, which emits \textbf{three} distinct vectors for 600
images with 92.8\% of them identical---the same collapse the relay model shows on PadChest
(three vectors, 99.6\%). Unresponsiveness of this size is therefore not a capacity or a
serving artefact; it tracks the modality.

\paragraph{A positive null separation is not guaranteed.} Every dermoscopy cell has a
positive $\Delta_{\text{null}}$ (0.155--0.504), but the sign is model-dependent elsewhere:
Qwen3-VL on brain MRI gives $-0.059$ [$-0.112$, $-0.003$] and MedGemma-1.5 on retinal OCT
gives $-0.058$ [$-0.080$, $-0.038$]. The phase-scrambled blank on the same OCT cells gives
$+0.003$ [$-0.028$, $0.032$], i.e.\ indistinguishable from zero. The 12 $\Delta_{\text{null}}$
estimates in Table~\ref{tab:crossmodal4} form one family, so we also apply a Bonferroni
correction ($\alpha=0.05/12=0.0042$). Under it, the Qwen3-VL MRI interval becomes
[$-0.137$, $+0.021$] and the Qwen3-VL OCT interval becomes [$-0.013$, $+0.131$]: both
include zero. The MedGemma-1.5 OCT interval remains negative, [$-0.090$, $-0.027$].
P2 therefore fails outright in one model--modality cell after correction, and nominally
in a second; a positive separation on one blank kind does not certify the other. The relay
model is positive on all three modalities (0.155--0.390), which shows that the sign is a
property of the model--modality--blank triple, not of the modality alone.

\paragraph{Contract coverage travels with responsiveness.} P1 alone would rank
Qwen3-VL as the least responsive model on dermoscopy (0.184), but its replies name only
18.7\% of the labels, so the status vector is mostly \texttt{not\_mentioned} and the
statistic is a lower bound from a sparse contract rather than a like-for-like comparison
with InternVL3.5's 0.289 at completeness 1.000. MedGemma-1.5 sits at 0.678 completeness
because 11.1\% of its clean cells and 12.4\% of its contrast cells are truncated. The
reportable unit is the pair (P1, completeness); reported alone, P1 can make a sparse
responder look unresponsive and a collapsed responder look responsive.

\paragraph{The operating point moves in opposite directions across models.}
Table~\ref{tab:crossmodal-op4} holds the corruption fixed and varies the model. Under
noise-severe on dermoscopy, the positive-call rate on label-negative lesions rises by
13.4\,pp for MedGemma-1.5 (0.536 $\rightarrow$ 0.670) but falls by 18.8\,pp for
InternVL3.5 (0.630 $\rightarrow$ 0.442), 8.8\,pp for Qwen3-VL (0.500 $\rightarrow$ 0.412)
and 17.4\,pp for the relay model (0.488 $\rightarrow$ 0.314); recall moves in the same
direction (+0.8, $-25.5$, $-6.2$ and $-25.0$\,pp). A recall-only table would describe the
first model as robust and the other three as degraded, although all four changed their
operating point. On retinal OCT the panel is conservative or collapsed: InternVL3.5 never
calls a negative positive (0.000) and keeps recall 0.918, while MedGemma-1.5 calls 8.7\%
positive with recall 0.067---two opposite failure forms that a single accuracy number
cannot distinguish.

\begin{table}[t]
\centering
\small
\caption{Operating point on the two modalities with a clinical negative class. ``Positive
rate'' is the share of label-negative images called positive; ``recall'' is the share of
label-positive images called positive. Clean and noise-severe cells only; the negative
class is the complement of the three positive labels of each modality.}
\label{tab:crossmodal-op4}
\begin{tabular}{l l r r r r r}
\toprule
Dataset & Model & $n_{\text{neg}}$ & Pos.\ rate clean & Pos.\ rate noise & Recall clean & Recall noise \\
\midrule
Dermoscopy & MedGemma-1.5-4B & 330 & 0.536 & 0.670 & 0.810 & 0.818 \\
           & InternVL3.5-8B & 330 & 0.630 & 0.442 & 0.891 & 0.636 \\
           & Qwen3-VL-8B & 330 & 0.500 & 0.412 & 0.713 & 0.651 \\
           & DeepSeek-V4.1-Flash (relay) & 172 & 0.488 & 0.314 & 0.836 & 0.586 \\
\midrule
Retina OCT & MedGemma-1.5-4B & 150 & 0.087 & 0.067 & 0.067 & 0.029 \\
           & InternVL3.5-8B & 150 & 0.000 & 0.000 & 0.918 & 0.902 \\
           & Qwen3-VL-8B & 150 & 0.013 & 0.000 & 0.480 & 0.224 \\
           & DeepSeek-V4.1-Flash (relay) & 70 & 0.000 & 0.000 & 0.748 & 0.739 \\
\bottomrule
\end{tabular}
\end{table}

\paragraph{Limits of this panel.} Three local checkpoints and one relay endpoint, six cells
per image rather than the 22-cell grid: no resolution, blur, JPEG or brightness step, and a
single synthetic seed. OCT carries no patient identifier, so its intervals resample images
rather than patients; dermoscopy is resampled at the lesion level and MRI at the patient
level. No human reader is included. These limits bound the magnitudes, not the ordering:
the modality ordering replicates in all four models, and the MedGemma-1.5 OCT separation
is the only negative one that survives the 12-test correction; the Qwen3-VL MRI
separation is nominal.
